\begin{proof}[Proof of \cref{obj:lem:finite-private-polynomial-calibration}]
The resource assumptions give
\[
\delta=\frac{e^\varepsilon-1}{e^\varepsilon+1}>0,
\qquad b>0,
\qquad \sigma^2=\frac{b^2}{m}>0,
\qquad \sigma>0,
\qquad L\ge1.
\]
All statistics below are measurable and have finite moments, since their message spaces are finite.

\begin{enumerate}
\item \textit{Global polynomial bias.}
Fix an even degree \(D\ge2\) satisfying the first set of degree conditions. In particular, \(D\le m\). By the moment unbiasedness in \cref{obj:lem:private-moment-and-dependence}, for every positive radius \(a\),
\[
\mathbb E\widehat p_{a,D,j}
=\sum_{v=0}^D a_{Dv}a^{1-v}\tau_j(P)^v
=a q_D\!\left(\frac{\tau_j(P)}a\right)
=p_{a,D}(\tau_j(P)).
\]
The interior means in \cref{obj:ass:interior-means} imply
\[
-\frac12\le\tau_j(P)\le\frac12.
\]
Consequently, at \(a=1/2\), the assumed approximation bound gives
\[
\left|\mathbb E\widehat p_{1/2,D,j}-|\tau_j(P)|\right|
=\frac12\left|q_D(2\tau_j(P))-|2\tau_j(P)|\right|
\le\frac1{\pi(D+1)}
\le\frac1{D+1}.
\]
% lean: CausalSmith.Stat.LdpOptvalueUniformFrontier.global_polynomial_bias_of_gate

\item \textit{Global polynomial variance.}
For \(1\le v\le D\), the moment second-moment bound in \cref{obj:lem:private-moment-and-dependence} applies because \(2v\le m\). It yields
\[
\begin{aligned}
\mathbb E\left[\left((1/2)^{1-v}\mathsf U_{vj}\right)^2\right]
&=\frac14\,4^v\mathbb E[\mathsf U_{vj}^2]\\
&\le\frac14\left(4\tau_j(P)^2+8v\sigma^2\right)^v\\
&\le\frac14(1+8D\sigma^2)^D.
\end{aligned}
\]
For \(v=0\), the same final bound follows from \(\mathsf U_{0j}=1\). Here we used \(\tau_j(P)^2\le1/4\) and \(1+8D\sigma^2\ge1\).

The assumed coefficient bound and \(\log2\le1\) give, for \(0\le v\le D\),
\[
|a_{Dv}|
\le 2^{3D/2}
=\exp\!\left(\frac{3D}{2}\log2\right)
\le e^{3D/2},
\qquad
a_{Dv}^2\le e^{3D}.
\]
Finite Cauchy--Schwarz therefore gives
\[
\begin{aligned}
\mathbb E[\widehat p_{1/2,D,j}^{\,2}]
&\le(D+1)\sum_{v=0}^D
 a_{Dv}^2
 \mathbb E\left[\left((1/2)^{1-v}\mathsf U_{vj}\right)^2\right]\\
&\le\frac14(D+1)^2e^{3D}(1+8D\sigma^2)^D\\
&\le\frac14e^{5D}(1+8D\sigma^2)^D,
\end{aligned}
\]
where the last inequality uses
\[
D+1\le e^D,
\qquad (D+1)^2\le e^{2D}.
\]

For the stipulated \(H=\log(e+\sigma^2L)\), we have
\[
H\ge1,
\qquad
DH\le\frac L{1024},
\qquad D\le L,
\qquad e^H=e+\sigma^2L.
\]
Since \(8\le e^3\), these inequalities imply
\[
1+8D\sigma^2
\le8(e+\sigma^2L)
=8e^H
\le e^{4H}.
\]
Thus
\[
e^{5D}(1+8D\sigma^2)^D
\le e^{5D+4DH}
\le e^{9DH},
\]
and hence
\[
\operatorname{Var}(\widehat p_{1/2,D,j})
\le\mathbb E[\widehat p_{1/2,D,j}^{\,2}]
\le\frac14e^{9DH}.
\]
Each polynomial estimate is a deterministic measurable statistic of its own column. Applying the average-variance conclusion of \cref{obj:lem:private-moment-and-dependence}, with \(0<m\le n\), proves
\[
\operatorname{Var}\left(\frac1d\sum_{j=1}^d\widehat p_{1/2,D,j}\right)
\le\frac2d\max_{j\in[d]}
       \operatorname{Var}(\widehat p_{1/2,D,j})
\le\frac{e^{9DH}}{2d}.
\]
% lean: CausalSmith.Stat.LdpOptvalueUniformFrontier.global_polynomial_calibration_of_gate

\item \textit{Hybrid bias inside the approximation window.}
Now fix a degree satisfying the second set of hypotheses. The floor inequalities give
\[
D\ \text{is even},
\qquad
D>\frac L{1024}-2,
\qquad
2\le D,
\qquad
\frac L{2048}\le D\le\frac L{1024},
\]
where \(L\ge4096\) supplies the lower bounds. Also \(D\le m\), and
\[
T_*=4\sigma\sqrt L>0,
\qquad
a=2T_*=8\sigma\sqrt L>0,
\qquad
a^2=64\sigma^2L.
\]

The two release blocks are independent by \cref{obj:ass:iid-people} and the independent participant releases. Within either block, the rows of a fixed column are independent. The row moments in \cref{obj:lem:private-moment-and-dependence} therefore give
\[
\mathbb EM_j=\mathbb E\overline W_j=\tau_j(P),
\]
and
\[
\begin{aligned}
\mathbb E[\overline W_j^{\,2}]
&=\frac{mb^2+m(m-1)\tau_j(P)^2}{m^2}
=\tau_j(P)^2+\frac{b^2-\tau_j(P)^2}{m},\\
\operatorname{Var}(\overline W_j)
&=\frac{b^2-\tau_j(P)^2}{m}
\le\sigma^2.
\end{aligned}
\]

First suppose that \(\tau_j(P)\ge0\). Define the pilot probabilities
\[
h_j=\Pr\{|M_j|\le T_*\},
\qquad
w_j=\Pr\{M_j<-T_*\}.
\]
They satisfy
\[
0\le h_j\le1,
\qquad w_j\ge0,
\qquad
h_j+\Pr\{M_j>T_*\}+w_j=1.
\]
The bounded mean concentration assumption applies to the independent pilot messages, each bounded by \(b\) in absolute value. For every real \(u\ge0\), it gives
\[
\Pr\{M_j-\tau_j(P)\ge u\}
\le e^{-u^2/(2\sigma^2)},
\qquad
\Pr\{M_j-\tau_j(P)\le-u\}
\le e^{-u^2/(2\sigma^2)}.
\]
In particular,
\[
\begin{aligned}
w_j
&\le\exp\!\left(-\frac{(\tau_j(P)+T_*)^2}{2\sigma^2}\right)\\
&\le\exp\!\left(-\frac{\tau_j(P)^2}{2\sigma^2}\right)e^{-8L},
\end{aligned}
\]
because \(T_*^2=16\sigma^2L\) and \(\tau_j(P)T_*\ge0\).

The elementary exponential envelopes give
\[
\tau_j(P)e^{-\tau_j(P)^2/(2\sigma^2)}\le\sigma,
\qquad
\tau_j(P)^2e^{-\tau_j(P)^2/(2\sigma^2)}\le\sigma^2.
\]
For the first inequality, multiply
\[
\frac{\tau_j(P)}{\sigma}
\le1+\frac{\tau_j(P)^2}{2\sigma^2}
\le e^{\tau_j(P)^2/(2\sigma^2)}
\]
by \(\sigma e^{-\tau_j(P)^2/(2\sigma^2)}\). For the second, use the elementary maximum
\[
u e^{-u}\le e^{-1}\qquad(u\ge0)
\]
at \(u=\tau_j(P)^2/(2\sigma^2)\), together with \(2/e\le1\). Consequently,
\[
\tau_j(P)w_j\le\sigma e^{-8L},
\qquad
\tau_j(P)^2w_j\le\sigma^2e^{-8L}\le\sigma^2.
\]
Moreover,
\[
\sqrt L\le L\le e^{8L},
\qquad
e^{-8L}\le\frac1{\sqrt L},
\qquad
\tau_j(P)w_j\le\frac{\sigma}{\sqrt L}.
\]

Suppose in addition that \(\tau_j(P)\le a\). Approximation and the degree lower bound imply
\[
\begin{aligned}
|p_{a,D}(\tau_j(P))-\tau_j(P)|
&\le\frac{2a}{\pi(D+1)}
\le\frac a{D+1}\\
&\le16384\,\frac{\sigma}{\sqrt L}.
\end{aligned}
\]
To compute the hybrid mean, its three disjoint pilot branches give
\[
H_j
=\widehat p_{a,D,j}\mathbf1_{\{|M_j|\le T_*\}}
+\overline W_j\mathbf1_{\{M_j>T_*\}}
-\overline W_j\mathbf1_{\{M_j<-T_*\}}.
\]
Block independence, polynomial unbiasedness, and the probability partition yield
\[
\begin{aligned}
\mathbb EH_j
&=h_jp_{a,D}(\tau_j(P))
 +(1-h_j-2w_j)\tau_j(P),\\
\mathbb EH_j-\tau_j(P)
&=h_j\bigl(p_{a,D}(\tau_j(P))-\tau_j(P)\bigr)
 -2\tau_j(P)w_j.
\end{aligned}
\]
Hence, throughout \(0\le\tau_j(P)\le a\),
\[
\begin{aligned}
|\mathbb EH_j-\tau_j(P)|
&\le h_j|p_{a,D}(\tau_j(P))-\tau_j(P)|
       +2\tau_j(P)w_j\\
&\le(16384+2)\frac{\sigma}{\sqrt L}
\le20000\,\frac{\sigma}{\sqrt L}.
\end{aligned}
\]
% lean: CausalSmith.Stat.LdpOptvalueUniformFrontier.hybrid_inside_nonneg_bias

\item \textit{Hybrid bias outside the approximation window.}
Consider the complementary case \(\tau_j(P)>a\), and define
\[
y_j=\frac{\tau_j(P)}a>1.
\]
On the central pilot event,
\[
M_j-\tau_j(P)
\le T_*-\tau_j(P)
\le-\frac{\tau_j(P)}2.
\]
The pilot concentration bound therefore gives
\[
h_j
\le\exp\!\left(-\frac{\tau_j(P)^2}{8\sigma^2}\right)
=e^{-8Ly_j^2}.
\]
The coefficient bound gives the polynomial growth estimate
\[
\begin{aligned}
|q_D(y_j)|
&\le\sum_{v=0}^D|a_{Dv}|y_j^v\\
&\le(D+1)e^{3D/2}y_j^D
\le e^{3D}y_j^D.
\end{aligned}
\]

For \(y_j\ge1\), the inequality \(\log y_j\le y_j-1\) and the nonnegativity of \((y_j-1)^2\) imply
\[
\log y_j\le\frac{y_j^2}{2}.
\]
For every nonnegative integer \(v\le L/1024\), it follows that
\[
3v+v\log y_j
\le\frac72v y_j^2
\le\frac7{2048}Ly_j^2
\le Ly_j^2,
\]
so
\[
e^{-8Ly_j^2}e^{3v}y_j^v\le e^{-7Ly_j^2}.
\]
Applying this with \(v=D\) proves
\[
h_j|p_{a,D}(\tau_j(P))|
\le a e^{-8Ly_j^2}e^{3D}y_j^D
\le a e^{-7Ly_j^2}.
\]
Since \(1,2\le L/1024\), the same bound with \(v=1,2\) gives
\[
h_jy_j^v
\le e^{-8Ly_j^2}y_j^v
\le e^{-8Ly_j^2}e^{3v}y_j^v
\le e^{-7Ly_j^2}
\qquad(v=1,2).
\]
Thus
\[
h_j\tau_j(P)\le a e^{-7Ly_j^2},
\qquad
h_j\tau_j(P)^2\le a^2e^{-7Ly_j^2}.
\]

The exponential lower bound \(e^{7L}\ge1+7L\ge L\) implies
\[
e^{-7Ly_j^2}\le e^{-7L}\le\frac1L,
\qquad
\frac aL=\frac{8\sigma}{\sqrt L}.
\]
Using the hybrid mean identity and the wrong-sign bound consequently gives
\[
\begin{aligned}
|\mathbb EH_j-\tau_j(P)|
&\le h_j|p_{a,D}(\tau_j(P))|
      +h_j\tau_j(P)+2\tau_j(P)w_j\\
&\le(8+8+2)\frac{\sigma}{\sqrt L}
\le20000\,\frac{\sigma}{\sqrt L}.
\end{aligned}
\]
Together with the inside case, this proves the bias bound for every nonnegative contrast, including zero and the boundary \(\tau_j(P)=a\).
% lean: CausalSmith.Stat.LdpOptvalueUniformFrontier.hybrid_nonneg_bias

\item \textit{Reflection for the hybrid bias.}
For a coordinate with \(\tau_j(P)<0\), define the reflected causal law
\[
P^{\mathrm{ref}}=G_{-\tau(P)},
\]
using \cref{obj:def:symmetric-law}. The parameter vector belongs to \([-1/2,1/2]^d\). The construction gives a probability law with uniform strata, fair assignment, and consistency, and
\[
\mu_{1j}(P^{\mathrm{ref}})=\frac{1-\tau_j(P)}2,
\qquad
\mu_{0j}(P^{\mathrm{ref}})=\frac{1+\tau_j(P)}2.
\]
Both means lie in \([1/4,3/4]\). Thus
\[
P^{\mathrm{ref}}\in\mathcal V_d,
\qquad
\tau_j(P^{\mathrm{ref}})=-\tau_j(P)
\quad(j\in[d]).
\]

We establish the reflection identity for the hybrid mean. The Chebyshev polynomials appearing in the explicit approximation are characterized, for \(\zeta\in\mathbb R\), by
\[
\mathcal T_0(\zeta)=1,
\qquad
\mathcal T_1(\zeta)=\zeta,
\qquad
\mathcal T_{v+1}(\zeta)
=2\zeta\mathcal T_v(\zeta)-\mathcal T_{v-1}(\zeta)
\quad(v\ge1).
\]
Induction gives
\[
\mathcal T_v(-\zeta)=(-1)^v\mathcal T_v(\zeta).
\]
The explicit approximation has the expansion
\[
q_D(\zeta)
=\frac2\pi+\frac4\pi
 \sum_{v=1}^{\lfloor D/2\rfloor}
 \frac{(-1)^{v+1}}{4v^2-1}\mathcal T_{2v}(\zeta).
\]
Consequently,
\[
q_D(-\zeta)=q_D(\zeta),
\qquad
(-1)^v a_{Dv}=a_{Dv}\quad(0\le v\le D),
\]
where the coefficient identity follows by comparing the two polynomials.

Viewing moments and estimates as functions of their message arrays, negation gives
\[
\mathsf U_{vj}(-W)=(-1)^v\mathsf U_{vj}(W),
\qquad
\widehat p_{a,D,j}(-W^{\mathrm{ev}})
=\widehat p_{a,D,j}(W^{\mathrm{ev}}).
\]
The pilot and evaluation means change sign, while
\[
|-M_j|=|M_j|,
\qquad
\operatorname{sign}(-M_j)(-\overline W_j)
=\operatorname{sign}(M_j)\overline W_j.
\]
These identities include \(M_j=0\) under the stipulated sign convention. They prove
\[
H_j(-W^{\mathrm{pil}},-W^{\mathrm{ev}})
=H_j(W^{\mathrm{pil}},W^{\mathrm{ev}}),
\]
with threshold ties preserved.

It remains to identify the reflected message law. Fair assignment, consistency, and uniform strata, as specified in
\cref{obj:ass:fair-randomization,obj:ass:consistency,obj:ass:uniform-covariates}, give
\[
\begin{aligned}
\mathbb E_P[S\mathbf1_{\{X=j\}}]
&=\frac12\left(2P(X=j,Y^1=1)-P(X=j)\right)\\
&\quad-\frac12\left(2P(X=j,Y^0=1)-P(X=j)\right)\\
&=\frac{\tau_j(P)}d.
\end{aligned}
\]
Averaging the vector kernel over a participant record therefore gives, for every row \(i\) in either block and every \(z\in\{-1,1\}^d\),
\[
\begin{aligned}
\Pr_P\!\left\{\left(\frac{W_{ij}}b\right)_{j=1}^d=z\right\}
&=2^{-d}\left(1+\delta\sum_{j=1}^d
 z_j\mathbb E_P[S\mathbf1_{\{X=j\}}]\right)\\
&=2^{-d}\left(1+\frac\delta d\sum_{j=1}^d\tau_j(P)z_j\right).
\end{aligned}
\]
Replacing \(z\) by \(-z\) gives precisely the row law under \(P^{\mathrm{ref}}\). Independence of participant rows and of the two blocks then yields
\[
\mathcal L_{P^{\mathrm{ref}}}(W^{\mathrm{pil}},W^{\mathrm{ev}})
=\mathcal L_P(-W^{\mathrm{pil}},-W^{\mathrm{ev}}).
\]
Combining this law identity with the invariance of \(H_j\) proves
\[
\mathbb E_{P^{\mathrm{ref}}}H_j=\mathbb E_PH_j.
\]
The nonnegative-contrast bias bound applied under \(P^{\mathrm{ref}}\) now gives
\[
\begin{aligned}
\left|\mathbb E_PH_j-|\tau_j(P)|\right|
&=\left|\mathbb E_{P^{\mathrm{ref}}}H_j
       -\tau_j(P^{\mathrm{ref}})\right|\\
&\le20000\,\frac{\sigma}{\sqrt L}.
\end{aligned}
\]
This proves the coordinate bias bound for every \(j\).
% lean: CausalSmith.Stat.LdpOptvalueUniformFrontier.hybrid_bias_bound

\item \textit{Polynomial second moments at the hybrid radius.}
For \(1\le v\le D\), \cref{obj:lem:private-moment-and-dependence} and \(2D\le m\) give
\[
\begin{aligned}
\mathbb E[(a^{1-v}\mathsf U_{vj})^2]
&\le a^2
 \left(\frac{\tau_j(P)^2+2v\sigma^2}{a^2}\right)^v\\
&\le a^2
 \left(\max\left\{1,\frac{\tau_j(P)^2+2D\sigma^2}{a^2}\right\}\right)^D.
\end{aligned}
\]
For \(v=0\), the left side is \(a^2\), so the final bound also holds. Finite Cauchy--Schwarz, \(a_{Dv}^2\le e^{3D}\), and \((D+1)^2\le e^{2D}\) yield
\[
\begin{aligned}
\mathbb E[\widehat p_{a,D,j}^{\,2}]
&\le(D+1)\sum_{v=0}^D
 a_{Dv}^2\mathbb E[(a^{1-v}\mathsf U_{vj})^2]\\
&\le a^2(D+1)^2e^{3D}
 \left(\max\left\{1,\frac{\tau_j(P)^2+2D\sigma^2}{a^2}\right\}\right)^D\\
&\le a^2e^{5D}
 \left(\max\left\{1,\frac{\tau_j(P)^2+2D\sigma^2}{a^2}\right\}\right)^D.
\end{aligned}
\]
% lean: CausalSmith.Stat.LdpOptvalueUniformFrontier.radius_polynomial_second_bound

\item \textit{Hybrid variance inside the approximation window.}
First suppose \(0\le\tau_j(P)\le a\). Since
\[
2D\sigma^2\le64\sigma^2L=a^2,
\]
we have
\[
\max\left\{1,\frac{\tau_j(P)^2+2D\sigma^2}{a^2}\right\}\le2.
\]
Using \(2\le e\) and \(D\le L/1024\), the polynomial second-moment bound becomes
\[
\mathbb E[\widehat p_{a,D,j}^{\,2}]
\le a^2e^{5D}2^D
\le a^2e^{6D}
\le a^2e^{L/100}.
\]

The three disjoint pilot branches give the pointwise identity
\[
\begin{aligned}
(H_j-\tau_j(P))^2
&=\mathbf1_{\{|M_j|\le T_*\}}
       (\widehat p_{a,D,j}-\tau_j(P))^2\\
&\quad+\mathbf1_{\{M_j>T_*\}}
       (\overline W_j-\tau_j(P))^2\\
&\quad+\mathbf1_{\{M_j<-T_*\}}
       (-\overline W_j-\tau_j(P))^2.
\end{aligned}
\]
Unbiasedness of the evaluation mean implies
\[
\begin{aligned}
\mathbb E[(-\overline W_j-\tau_j(P))^2]
&=\mathbb E[(\overline W_j-\tau_j(P))^2]
   +4\tau_j(P)\mathbb E\overline W_j\\
&=\operatorname{Var}(\overline W_j)+4\tau_j(P)^2.
\end{aligned}
\]
Block independence and the pilot probability partition therefore yield
\[
\begin{aligned}
\mathbb E[(H_j-\tau_j(P))^2]
&=h_j\mathbb E[(\widehat p_{a,D,j}-\tau_j(P))^2]\\
&\quad+(1-h_j)\operatorname{Var}(\overline W_j)
       +4\tau_j(P)^2w_j.
\end{aligned}
\]
The pointwise inequality
\[
(\widehat p_{a,D,j}-\tau_j(P))^2
\le2\widehat p_{a,D,j}^{\,2}+2\tau_j(P)^2
\]
and \(\operatorname{Var}(\overline W_j)\le\sigma^2\) give
\[
\begin{aligned}
\operatorname{Var}(H_j)
&\le\mathbb E[(H_j-\tau_j(P))^2]\\
&\le2h_j\mathbb E[\widehat p_{a,D,j}^{\,2}]
     +2h_j\tau_j(P)^2+\sigma^2+4\tau_j(P)^2w_j.
\end{aligned}
\]
This centered-error bound also applies to the outside case.

Inside the window, \(h_j\le1\), \(\tau_j(P)^2\le a^2\), and the wrong-sign estimate gives \(\tau_j(P)^2w_j\le\sigma^2\). Hence
\[
\begin{aligned}
\operatorname{Var}(H_j)
&\le2a^2e^{L/100}+2a^2+5\sigma^2\\
&=128\sigma^2Le^{L/100}+128\sigma^2L+5\sigma^2\\
&\le400\sigma^2Le^{L/100},
\end{aligned}
\]
using \(L\ge1\) and \(e^{L/100}\ge1\).
% lean: CausalSmith.Stat.LdpOptvalueUniformFrontier.hybrid_inside_nonneg_variance

\item \textit{Hybrid variance outside the approximation window.}
Suppose \(\tau_j(P)>a\), and again use \(y_j=\tau_j(P)/a>1\). Then
\[
\frac{\tau_j(P)^2+2D\sigma^2}{a^2}
=y_j^2+\frac{2D\sigma^2}{a^2}
\le y_j^2+1
\le2y_j^2,
\]
so
\[
\max\left\{1,\frac{\tau_j(P)^2+2D\sigma^2}{a^2}\right\}
\le2y_j^2.
\]
Combining the pilot tail and polynomial second-moment bounds gives
\[
h_j\mathbb E[\widehat p_{a,D,j}^{\,2}]
\le a^2e^{-8Ly_j^2}e^{5D}(2y_j^2)^D.
\]
Since \(2\le e\),
\[
(2y_j^2)^D\le e^D e^{2D\log y_j}.
\]
Moreover, \(\log y_j\le y_j^2/2\), \(y_j^2\ge1\), and \(D\le L/1024\) imply
\[
6D+2D\log y_j
\le7Dy_j^2
\le\frac7{1024}Ly_j^2
\le Ly_j^2.
\]
It follows that
\[
h_j\mathbb E[\widehat p_{a,D,j}^{\,2}]
\le a^2e^{-7Ly_j^2}\le a^2.
\]
The contrast-weighted pilot bound also gives
\[
h_j\tau_j(P)^2\le a^2e^{-7Ly_j^2}\le a^2.
\]
Using these inequalities and \(\tau_j(P)^2w_j\le\sigma^2\) in the centered-error bound yields
\[
\begin{aligned}
\operatorname{Var}(H_j)
&\le4a^2+5\sigma^2\\
&=256\sigma^2L+5\sigma^2\\
&\le400\sigma^2Le^{L/100}.
\end{aligned}
\]
The inside and outside cases establish the variance bound for every nonnegative contrast.
% lean: CausalSmith.Stat.LdpOptvalueUniformFrontier.hybrid_nonneg_variance

\item \textit{Reflection for the hybrid variance.}
If \(\tau_j(P)<0\), use the same law \(P^{\mathrm{ref}}=G_{-\tau(P)}\). The reflected message-law identity and invariance of \(H_j\), together with equality of its means, imply
\[
\begin{aligned}
\operatorname{Var}_{P^{\mathrm{ref}}}(H_j)
&=\mathbb E_{P^{\mathrm{ref}}}
       [(H_j-\mathbb E_{P^{\mathrm{ref}}}H_j)^2]\\
&=\mathbb E_P[(H_j-\mathbb E_PH_j)^2]
=\operatorname{Var}_P(H_j).
\end{aligned}
\]
Since \(\tau_j(P^{\mathrm{ref}})>0\), the nonnegative-contrast variance bound gives
\[
\operatorname{Var}_P(H_j)
=\operatorname{Var}_{P^{\mathrm{ref}}}(H_j)
\le400\sigma^2Le^{L/100}.
\]
Thus the asserted coordinate variance bound holds for every \(j\).
% lean: CausalSmith.Stat.LdpOptvalueUniformFrontier.hybrid_variance_bound

\item \textit{Variance of the hybrid average.}
Stack the pilot rows followed by the evaluation rows. Their joint law is the \(2m\)-participant product release law, and each \(H_j\) is a deterministic measurable function of the combined column
\[
(W^{\mathrm{pil}}_{1j},\ldots,W^{\mathrm{pil}}_{mj},
 W^{\mathrm{ev}}_{1j},\ldots,W^{\mathrm{ev}}_{mj})
\in\mathbb R^{2m}.
\]
Because \(0<2m\le n\), \cref{obj:lem:private-moment-and-dependence} applies to this combined block. The coordinate variance bounds therefore give
\[
\operatorname{Var}\left(\frac1d\sum_{j=1}^dH_j\right)
\le\frac2d\max_{j\in[d]}\operatorname{Var}(H_j)
\le\frac{800\sigma^2Le^{L/100}}d.
\]

We absorb the dimension factor using \(L\ge4096\). The elementary exponential lower bound gives
\[
\frac L4\le e^{L/4},
\qquad
L^2\le16e^{L/2},
\]
and
\[
4\le\frac{49L}{100}
=1+\left(\frac{49L}{100}-1\right)
\le e^{49L/100-1}.
\]
Since \(e^{L-1}=d\), it follows that
\[
\begin{aligned}
L^2e^{L/100}
&\le16e^{L/2}e^{L/100}\\
&=16e^{51L/100}\\
&\le4e^{51L/100}e^{49L/100-1}
=4d.
\end{aligned}
\]
As \(d,L>0\), this proves
\[
\frac{800\sigma^2Le^{L/100}}d
\le\frac{3200\sigma^2}{L},
\qquad
\operatorname{Var}\left(\frac1d\sum_{j=1}^dH_j\right)
\le\frac{3200\sigma^2}{L}.
\]
% lean: CausalSmith.Stat.LdpOptvalueUniformFrontier.hybrid_average_variance_envelope_le

\item \textit{Mean squared error of the hybrid average.}
Linearity of expectation, the definition of \(F\), and the coordinate bias bounds give
\[
\begin{aligned}
\left|\mathbb E\left[\frac1d\sum_{j=1}^dH_j\right]
      -F(\tau(P))\right|
&=\left|\frac1d\sum_{j=1}^d
       \bigl(\mathbb EH_j-|\tau_j(P)|\bigr)\right|\\
&\le\frac1d\sum_{j=1}^d
       \left|\mathbb EH_j-|\tau_j(P)|\right|\\
&\le20000\,\frac{\sigma}{\sqrt L}.
\end{aligned}
\]
Squaring, with \(\sigma^2>0\) and \(L>0\), yields
\[
\left(\mathbb E\left[\frac1d\sum_{j=1}^dH_j\right]
      -F(\tau(P))\right)^2
\le400000000\,\frac{\sigma^2}{L}.
\]
Finally, expansion around the expectation gives the variance-plus-squared-bias identity
\[
\begin{aligned}
\mathbb E\left[
 \left(\frac1d\sum_{j=1}^dH_j-F(\tau(P))\right)^2\right]
&=\operatorname{Var}\left(\frac1d\sum_{j=1}^dH_j\right)\\
&\quad+\left(\mathbb E\left[\frac1d\sum_{j=1}^dH_j\right]
             -F(\tau(P))\right)^2\\
&\le(3200+400000000)\frac{\sigma^2}{L}\\
&=400003200\,\frac{\sigma^2}{L}\\
&\le500000000\,\frac{\sigma^2}{L}.
\end{aligned}
\]
% lean: CausalSmith.Stat.LdpOptvalueUniformFrontier.hybrid_average_mse_bound
\end{enumerate}
% lean: CausalSmith.Stat.LdpOptvalueUniformFrontier.finite_private_polynomial_calibration
\end{proof}
